\begin{question}
Recall the atomic broadcast problem from Exercise~\ref{ex:atomicBroadcast}: a
mechanism is required to allow a single sender to synchronously broadcast a
message to~|n| receivers, repeatedly.  Implement a class to solve this
problem, using conditions internally.  You may assume $\sm n > 0$.
\end{question}

%%%%%

\begin{answerI}
My code is below.  As with the barrier synchronisation example in the body of
the chapter, we need to ensure that different synchronisations do not
interfere.  To that end, all the receivers wait until the previous
synchronisation has finished and all threads have returned; they then record
their presence and wait for the sender.  The sender waits until all |n|
receivers have entered, writes its values, and signals to the receivers.  The
receivers pick up the sender's value; the last one to leave resets for the
next round and signals to waiting receivers.  Comments in the code give more
details. 
%
\begin{scala}
class AtomicBroadcast[A](n: Int) extends tacp.clientServer.AtomicBroadcastT[A]{
  require(n > 0)

  /** Number of receivers waiting (when £slotFilled = false£), or the number still
    * to leave (when £slotFilled = true£). */
  private var waiting = 0

  private var slot: A = _

  /** Is the value in £slot£ valid?  This remains true until the last receiver has
    * read £slot£. */
  private var slotFilled = false

  private val lock = new Lock

  /** A condition to signal to the sender that all the readers are waiting, and
    * so it can write its value and return. */
  private val canWrite = lock.newCondition

  /** A condition to signal to receivers that the previous synchronisation is
    * complete, so they can enter the main part of the operation. */
  private val canEnter = lock.newCondition

  /** A condition to signal to receivers that they can read the sender's value. */
  private val canRead = lock.newCondition

  def send(x: A) = lock.mutex{
    canWrite.await(slotFilled == false && waiting == n) 
                                             // Wait for the receivers (2).
    slot = x; slotFilled = true
    canRead.signalAll() // Signal to the receivers at (3).
  }

  def receive(): A = lock.mutex{
    canEnter.await(slotFilled == false) // Wait for previous round to finish (1).
    waiting += 1
    if(waiting == n) 
      canWrite.signal() // The last receiver signals to the sender at (2).
    canRead.await() // Wait for the sender (3).
    waiting -= 1
    if(waiting == 0){ // The last receiver to leave resets for the next round.
      slotFilled = false 
      canEnter.signalAll() // Signal to receivers on the next round at (1).
    }
    slot
  }
}
\end{scala}
%
This requires $\sm n > 0$, or else the sender will deadlock on its second
iteration: |slotFilled| will be |true|, and there is no receiver to set it to
|false|.  We could avoid this assumption by treating $\sm n = 0$ as a special
case.
\end{answerI}
